% series_pendant_reduction.tex —— 系统模型级串联与悬垂降阶

\section{串联化简与悬垂负荷折叠}\label{sec:gr-series-pendant}

\subsection{共同执行语义}

两个 reducer 都从 graph/system 副本开始，仅消费匹配类型的 plan actions。被消去 bus/branch 不从
内部结果中立即删除，而是标 \fld{in\_service=false}；HTTP 导出层另行 compact。动作大小不符、
母线/边找不到或边已停运时跳过。plan 中的边引用明确为生成该 plan 时的
\texttt{graph.edges[]} position；调用方不能把稳定组件 \fld{index} 填入该字段。

\subsection{串联适用条件}

计划层已要求被消去节点无注入、无 shunt/control、在运度数 2，且两条边均为普通 AC/DC line。
若保留热限，则两边必须没有 rate\_a\_mva。执行层不重复检查这些工程条件，信任 plan；手工构造
ReductionPlan 的调用方承担前置条件责任。

对 $i-j-k$，代码构造
\begin{align}
r_{ik}&=r_{ij}+r_{jk},&x_{ik}&=x_{ij}+x_{jk},\\
b_{ik}&=b_{ij}+b_{jk},&tap_{ik}&=1,&shift_{ik}&=0,\\
rate_{ik}&=0.
\end{align}
串联阻抗相加对无注入节点的 series impedance 精确；两个 $\pi$ 型支路的总充电电纳直接相加不是
一般严格二端口级联等值，且原热限被清零。故 result 的“exact”必须读作
\textbf{零注入串联阻抗关系精确}，不是所有支路参数/约束逐项保真。

\subsection{新增边与 ID}

新 branch component index 从
\begin{equation}
\max\{\,1+\max ACBranch.index,\ 1+\max DCBranch.index\,\}
\end{equation}
开始，并赋给 new edge 的 \fld{comp\_index} 与新 AC/DCBranch.index。\fld{edge\_id} 则始终取
push-back 前的 \texttt{reduced\_graph.edges.size()}，所以初始、收缩及 series 后都保持
\begin{equation}
\fld{edge\_id}=\texttt{graph.edges[]}\text{ position}.
\end{equation}
稳定组件 ID 与图存储位置是两个独立空间；对外支路证书使用带域 \texttt{BranchRef}。

\subsection{串联模型更新}

记录 \texttt{SeriesReductionRecord}：eliminated/from/to bus 为稳定 ID，domain 域限定，
original\_branch\_ids 存两条\textbf{源组件 comp\_index}，new\_branch\_id 存新 branch component index。
旧 AC/DCBranch 按 comp\_index 停运，被消去 bus 停运，新等值 branch 加入对应域。

action 输入使用 graph edge position；mapping 与 record 始终使用稳定 branch component index。
\srcpath{set\_bus\_reduction/set\_branch\_reduction} 每次同步重建域限定反向表，因此 eliminated bus
会进入代表 bus 的 original members，两条源支路也会反向归因到新等值支路。

\subsection{多动作边界}

每个 action 从持续变更的 reduced graph 查边，但 plan 来自原图。相邻链节点在计划层被冲突集合
跳过，避免第二个动作引用已失效邻点；非相邻动作可一次执行。reduced graph 的原邻接条目不删除，
依靠旧 edge.in\_service=false 过滤。

\subsection{悬垂负荷聚合}

对 parent $i$ 与 leaf $j$，执行层从 bus-level demand 加上在运详细 Load/DCLoad。基准功率取
system.base\_mva，非正时回退 100 MVA。以 $|V_i|^2\approx1$ 计算
\begin{align}
p_j^{pu}&=P_j/S_B,&q_j^{pu}&=Q_j/S_B,\\
\ell_j&=(p_j^{pu})^2+(q_j^{pu})^2,\\
p_{abs}^{pu}&=p_j^{pu}+r_{ij}\ell_j,\\
q_{abs}^{pu}&=q_j^{pu}+x_{ij}\ell_j \quad(AC),
\end{align}
DC 固定 $q=0$。吸收量乘 $S_B$ 后写入 parent 的 bus-level demand；原详细 load 行仍留在已停运
leaf bus 下并标为停运，HTTP compact 会删除该停运行；其有功/无功已计入 parent 的 bus-level demand。

\subsection{悬垂记录与分支空间}

\texttt{PendantReductionRecord::p\_load\_absorbed/q\_load\_absorbed} 存 MW/MVAr。
\fld{source\_branch\_index} 明确存源 ACBranch/DCBranch 的稳定 \fld{index}；恢复按同一域和同一
组件空间查阻抗，找不到时抛 \texttt{invalid\_argument}，不再用 $Z=0$ 伪造结果。

Pendant result 从身份 mapping 开始，把 eliminated bus 正向映射到 parent，并把源 branch 映射到
invalid \texttt{BranchRef} 表示“已消去且无替代支路”；域限定正反向表随之重建。

\subsection{悬垂近似误差}

真实支路损耗为 $Z|S/V_j|^2$，而折叠使用 $|V_i|=1$。电压偏低、重载、高 $R/X$ 时误差增大；
折叠后的 PF 只证明近似系统可解，不证明原系统精确等价。当前结果不报告误差界、基准/化简损耗差、
近似有效标志或折叠恢复收敛状态。

\subsection{调用示例的正确顺序}

\begin{enumerate}
  \item build graph；
  \item classify + make plan；
  \item apply series；
  \item 在 series result 的 graph/system 上重新 classify + plan；
  \item 可选 apply pendant；
  \item 求解 reduced system；
  \item 按应用的逆序调用 pendant/series/switch recovery，并自行维护每阶段 mapping。
\end{enumerate}
各阶段 mapping 可用 \srcpath{compose\_reduction\_mappings} 按执行顺序合成一份证书；operation
records 同样按正向执行顺序串接，恢复仍应按阶段逆序调用相应恢复函数。
